\paragraph{Conjugate-coordinate class: where the flux eigenvalues of the Euclidean symbol coincide (conditional; bounded no-go).}
Keep the setting of the preceding paragraph: Poisson limit, leaf \(\mathbb{CP}^{1,2}\), \(D^a=\{T^a,\cdot\}\), the nilpotent (Landau) leading symbol with centre variable \(\zeta\in\mathbb R^2\), the contraction \(\eta\to\delta\) of the four matrix indices with the Lie--Poisson algebra \(\mathfrak{so}(4,2)\) kept, and the sheet realised by the orbit. Write \(y^a=M^{a5}\), \(y^4=M^{45}\), \(y\cdot y=-\ell^2\), and \(Q^{ab}=(y^4)^2\eta^{ab}+y^ay^b\), the Gram matrix of the fields \(D^a\) in the ambient metric of the hyperboloid; \(Q\) is positive definite with \(\det Q=\ell^2|y^4|^6\). At \(300\) random orbit points, with both centre directions, the matrix \(Q^{-1/2}F_{ab}Q^{-1/2}\), \(F_{ab}=\zeta_kS^{ab}{}_k\), is self-dual or anti-self-dual with respect to \(\delta\) to a relative accuracy of \(10^{-9}\), with the sign equal to \(\operatorname{sign}y^4\); this is the expected statement that the vertical curvature of the horizontal distribution over the constant-curvature base is self-dual for the ambient metric. Hence \(F=Q^{1/2}\omega Q^{1/2}\) with \(\omega\) self-dual, and the plane strengths of \(F\) are \(f_{1,2}=|p|\pm|r|\), with \(p\) and \(r\) the \(\delta\)-self-dual and anti-self-dual parts of \(F\). Let \(a\ge b\) be the eigenvalues of \(Q\) on the plane spanned by the time axis and \(\vec y\), so that \(ab=\ell^2(y^4)^2\) and the remaining eigenvalue is \((y^4)^2\) twice. For \(\zeta\) such that \(\omega\) points along the axis of that plane one finds \(f_2/f_1=\min(\ell,|y^4|)/\max(\ell,|y^4|)\); for \(\omega\) orthogonal to it, \(f_2/f_1=\sqrt{b/a}\); for other directions the ratio lies between these values. The numerical range of \(f_2/f_1\) at the sampled points lies inside this bracket to \(10^{-13}\).

\emph{Consequences.} (i) The vector sector has no negative mode at fixed \(\zeta\) only if \(f_1=f_2\), and a stable point must have this for every \(\zeta\), because the set of \(\zeta\) with \(f_1=f_2\) is a union of at most two lines unless \(F\) is (anti-)self-dual identically. The closed forms give \(f_1=f_2\) for all \(\zeta\) exactly on \(\vec y=0\), \(|y^4|=\ell\), that is on the two-sphere fibres over the points \(y^0=-\sqrt2\,\ell\), \(y^4=\pm\ell\): codimension \(4\) in the six-dimensional leaf (the minimisation of the distance to this set found it, with \(f_2/f_1=1\) to six digits at six points). On the rest of the leaf the lowest vector eigenvalue is \(-|f_1-f_2|<0\) for generic \(\zeta\), and \(|f_1-f_2|/(f_1+f_2)=(\ell-|y^4|)/(\ell+|y^4|)\) on the time axis grows linearly with the distance from the set. (ii) The same holds for the compact real form \(\mathfrak{so}(6)\) with \(\delta\) throughout: there \(f_2/f_1=|y^4|/\ell\) for every \(\zeta\) and every fibre point (checked at \(200\) points, error \(4\cdot10^{-16}\)), equal to \(1\) only over the poles \(\vec y=0\). A compact Euclidean section with the contraction \(\delta_{ab}\) therefore does not replace the instability by zero modes; self-duality holds identically only for the contraction by the position-dependent ambient metric \(Q^{-1}\), which is not a constant contraction of the matrix indices. (iii) Constant additional fluxes \(f_3,\dots\) from transverse matrices change the lowest level to \(\sum_kf_k-2\max_kf_k\) (Troost's levels \(2\pi(\sum_i(2m_i+1)f_i\pm2f_k)\), hep-th/9909187, eq.~(35), and the no-tachyon condition of its Section~9, \(f_1\le f_2+\dots+f_n\)); because \(f_1-f_2\) grows linearly in \(|\zeta|\), a \(\zeta\)-independent transverse flux restores stability only for \(|\zeta|\) below a bound set by that flux, that is on a finite fuzzy fibre, not in the Poisson limit. (iv) At the self-dual points the vector density \(\prod_j f_j/(4\pi\sinh sf_j)\) times the Zeeman trace tends to \(f^2/(2\pi^2)\) for \(sf\gg1\), independent of \(s\), so the integral over \(\zeta\) diverges there as well; the proper-time cutoff does not regulate it.

This is a statement for the Poisson limit, the leading symbol, equiregular points, the contraction \(\eta\to\delta\) (or \(\delta\) with \(\mathfrak{so}(6)\)), constant additional fluxes and bosonic content. It does not treat a contraction by a position-dependent metric, Moyal terms, fermions, a finite fibre, or other backgrounds for the transverse matrices, and it determines neither the sign nor the size of \(w\). This changes no ledger label; Problem~F remains open.
